\Chapter{33}

\Verse{1}
{
    \zA{却说王夫人唤上金钏儿的母亲来，拿了几件簪环当面赏了，又吩咐：“请几众 僧人念经超度他。” 金钏儿的母亲磕了头，谢了出去。
        原来宝玉会过雨村回来，听见金钏儿含羞自尽，心中早已五内摧伤，进来又被 王夫人数说教训了一番，也无可回说。 看见宝钗进来，方得便走出，茫然不知何往，
        背着手，低着头，一面感叹，一面慢慢的信步走至厅上。 刚转过屏门，不想对面来 了一人正往里走，可巧撞了个满怀。 只听那人喝一声：“站住！”}
}
{
    \eA{Madame Wang, for we shall now continue our story, sent for Chin Ch'uan-erh's
        mother. On her arrival, she gave her several hair-pins and rings, and then
        told her that she could invite several Buddhist priests as well to read the
        prayers necessary to release the spirit from purgatory. The mother
        prostrated herself and expressed her gratitude; after which, she took her
        leave.}
}
{
}

\Verse{2}
{
    \zA{宝玉唬了一跳， 抬头看时，不是别人，却是他父亲。 早不觉倒抽了一口凉气，只得垂手一旁站着。 贾政道：“好端端的，你垂头丧气的？ 什么？
        方才雨村来了要见你，那半天才出来! 既出来了，全无一点慷慨挥洒的谈吐，仍是委委琐琐的。 我看你脸上一团私欲愁闷
        气色!这会子又嗳声叹气，你那些还不足、还不自在？ 无故这样，是什么原故？” 宝 玉素日虽然口角伶俐，此时一心却为金钏儿感伤，恨不得也身亡命殒；}
}
{
    \eA{Indeed, Pao-yü, on his return from entertaining Yü-ts'un, heard the tidings
        that Chin Ch'uan-erh had been instigated by a sense of shame to take her own
        life and he at once fell a prey to grief. So much so, that, when he came
        inside, and was again spoken to and admonished by Madame Wang, he could not
        utter a single word in his justification.}
}
{
}

\Verse{3}
{
    \zA{如今见他父 亲说这些话，究竟不曾听明白了，只是怔怔的站着。 贾政见他惶悚，应对不似往日，原本无气的，这一来倒生了三分气。 方欲说话，
        忽有门上人来回：“忠顺亲王府里有人来，要见老爷。” 贾政听了，心下疑惑，暗 暗思忖道：“素日并不与忠顺府来往，为什么今日打发人来？” 一面想，一面命：
        “快请厅上坐。” 急忙进内更衣。 出来接见时，却是忠顺府长府官，一面彼此见了 礼，归坐献茶。}
}
{
    \eA{But as soon as he perceived Pao-ch'ai make her appearance in the room, he
        seized the opportunity to scamper out in precipitate haste. Whither he was
        trudging, he himself had not the least idea. But throwing his hands behind
        his back and drooping his head against his chest, he gave way to sighs,
        while with slow and listless step he turned towards the hall.}
}
{
}

\Verse{4}
{
    \zA{未及叙谈，那长府官先就说道：“下官此来，并非擅造潭府，皆因 奉命而来，有一件事相求。 看王爷面上，敢烦老先生做主，不但王爷知情，且连下
        官辈亦感谢不尽。” 贾政听了这话，摸不着头脑，忙陪笑起身问道：“大人既奉王 命而来，不知有何见谕？ 望大人宣明，学生好遵谕承办。” 那长府官冷笑道：“也
        不必承办，只用老先生一句话就完了。 我们府里有一个做小旦的琪官，一向好好在 府，如今竟三五日不见回去，各处去找，又摸不着他的道路。}
}
{
    \eA{Scarcely, however, had he rounded the screen-wall, which stood in front of
        the door-way, when, by a strange coincidence, he ran straight into the arms
        of some one, who was unawares approaching from the opposite direction, and
        was just about to go towards the inner portion of the compound. "Hallo!"
        that person was heard to cry out, as he stood still. Pao-yü sustained a
        dreadful start.}
}
{
}

\Verse{5}
{
    \zA{因此各处察访，这一 城内十停人倒有八停人都说：他近日和衔玉的那位令郎相与甚厚。 下官辈听了，尊 府不比别家，可以擅来索取，因此启明王爷。
        王爷亦说：‘若是别的戏子呢，一百 个也罢了； 只是这琪官，随机应答，谨慎老成，甚合我老人家的心境，断断少不得 此人。’
        故此求老先生转致令郎，请将琪官放回：一则可慰王爷谆谆奉恳之意，二 则下官辈也可免操劳求觅之苦。” 说毕，忙打一躬。
        贾政听了这话，又惊又气，即命唤宝玉出来。}
}
{
    \eA{Raising his face to see, he discovered that it was no other than his father.
        At once, he unconsciously drew a long breath and adopted the only safe
        course of dropping his arms against his body and standing on one side. "Why
        are you," exclaimed Chia Cheng, "drooping your head in such a melancholy
        mood, and indulging in all these moans?}
}
{
}

\Verse{6}
{
    \zA{宝玉也不知是何原故，忙忙赶来， 贾政便问：“该死的奴才!你在家不读书也罢了，怎么又做出这些无法无天的事来!
        那琪官现是忠顺王爷驾前承奉的人，你是何等草莽，无故引逗他出来，如今祸及于 我！” 宝玉听了，唬了一跳，忙回道：“实在不知此事。 究竟‘琪官’两个字，不
        知为何物，况更加以‘引逗’二字！” 说着便哭。 贾政未及开口，只见那长府官冷 笑道：“公子也不必隐饰。
        或藏在家，或知其下落，早说出来，我们也少受些辛苦， 岂不念公子之德呢！”}
}
{
    \eA{When Yü-ts'un came just now and he asked to see you, you only put in your
        appearance after a long while. But though you did come, you were not in the
        least disposed to chat with anything like cheerfulness and animation; you
        behaved, as you ever do, like a regular fool. I detected then in your
        countenance a certain expression of some hidden hankering and sadness; and
        now again here you are groaning and sighing!}
}
{
}

\Verse{7}
{
    \zA{宝玉连说：“实在不知。 恐是讹传，也未见得。” 那长府官 冷笑两声道：“现有证据，必定当着老大人说出来，公子岂不吃亏？ 既说不知，此
        人那红汗巾子怎得到了公子腰里？” 宝玉听了这话，不觉轰了魂魄，目瞪口呆。 心 下自思：“这话他如何知道？ 他既连这样机密事都知道了，大约别的瞒不过他。
        不 如打发他去了，免得再说出别的事来。” 因说道：“大人既知他的底细，如何连他 置买房舍这样大事倒不晓得了。}
}
{
    \eA{Does all you have not suffice to please you? Are you still dissatisfied?
        You've no reason to be like this, so why is it that you go on in this way?"}
}
{
}

\Verse{8}
{
    \zA{听得说他如今在东郊离城二十里有个什么紫檀堡， 他在那里置了几亩田地，几间房舍。 想是在那里，也未可知。” 那长府官听了，笑
        道：“这样说，一定是在那里了。 我且去找一回，若有了便罢； 若没有，还要来请 教。” 说着，便忙忙的告辞走了。
        贾政此时气得目瞪口歪，一面送那官员，一面回头命宝玉：“不许动!回来有 话问你！” 一直送那官去了。 才回身时，忽见贾环带着几个小厮一阵乱跑。 贾政喝
        命小厮：“给我快打！”}
}
{
    \eA{Pao-yü had ever, it is true, shown a glib tongue, but on the present
        occasion he was so deeply affected by Chin Ch'uan-erh's fate, and vexed at
        not being able to die that very instant and follow in her footsteps that
        although he was now fully conscious that his father was speaking to him he
        could not, in fact, lend him an ear, but simply stood in a timid and nervous
        mood.}
}
{
}

\Verse{9}
{
    \zA{贾环见了他父亲，吓得骨软筋酥，赶忙低头站住。 贾政便 问：“你跑什么？ 带着你的那些人都不管你，不知往那里去，由你野马一般！” 喝
        叫：“跟上学的人呢？” 贾环见他父亲甚怒，便乘机说道：“方才原不曾跑，只因 从那井边一过，那井里淹死了一个丫头，我看脑袋这么大，身子这么粗，泡的实在
        可怕，所以才赶着跑过来了。” 贾政听了，惊疑问道：“好端端，谁去跳井？ 我家 从无这样事情。}
}
{
    \eA{Chia Cheng noticed that he was in a state of trembling and fear, not as
        ready with an answer as he usually was, and his sorry plight somewhat
        incensed him, much though he had not at first borne him any ill-feeling. But
        just as he was about to chide him, a messenger approached and announced to
        him: "Some one has come from the mansion of the imperial Prince Chung Shun,
        and wishes to see you, Sir."}
}
{
}

\Verse{10}
{
    \zA{自祖宗以来，皆是宽柔待下，大约我近年于家务疏懒，自然执事人 操克夺之权，致使弄出这暴殒轻生的祸来。 若外人知道，祖宗的颜面何在！” 喝命：
        “叫贾琏、赖大来！” 小厮们答应了一声，方欲去叫，贾环忙上前拉住贾政袍襟， 贴膝跪下道：“老爷不用生气。 此事除太太屋里的人，别人一点也不知道。 我听见
        我母亲说――”说到这句，便回头四顾一看。 贾政知其意，将眼色一丢，小厮们明 白，都往两边后面退去。}
}
{
    \eA{At this announcement, surmises sprung up in Chia Cheng's mind. "Hitherto,"
        he secretly mused, "I've never had any dealings with the Chung Shun mansion,
        and why is it that some one is despatched here to-day?" As he gave way to
        these reflections. "Be quick," he shouted, "and ask him to take a seat in
        the pavilion," while he himself precipitately entered the inner room and
        changed his costume.}
}
{
}

\Verse{11}
{
    \zA{贾环便悄悄说道：“我母亲告诉我说：宝玉哥哥前日在太 太屋里，拉着太太的丫头金钏儿，强奸不遂，打了一顿，金钏儿便赌气投井死了。”
        话未说完，把个贾政气得面如金纸，大叫：“拿宝玉来！” 一面说，一面便往 书房去，喝命：“今日再有人来劝我，我把这冠带家私，一应就交与他和宝玉过去!
        我免不得做个罪人，把这几根烦恼鬓毛剃去，寻个干净去处自了，也免得上辱先人、 下生逆子之罪！”}
}
{
    \eA{When he came out to greet the visitor, he discovered that it was the senior
        officer of the Chung Shun mansion. After the exchange of the salutations
        prescribed by the rites, they sat down and tea was presented.}
}
{
}

\Verse{12}
{
    \zA{众门客仆从见贾政这个形景，便知又是为宝玉了，一个个咬指吐 舌，连忙退出。 贾政喘吁吁直挺挺的坐在椅子上，满面泪痕，一叠连声：“拿宝玉
        来!拿大棍拿绳来!把门都关上!有人传信到里头去，立刻打死！” 众小厮们只得齐 齐答应着，有几个来找宝玉。
        那宝玉听见贾政吩咐他“不许动”，早知凶多吉少，那里知道贾环又添了许多 的话？ 正在厅上旋转，怎得个人往里头捎信，偏偏的没个人来，连焙茗也不知在那 里。
        正盼望时，只见一个老妈妈出来。}
}
{
    \eA{But before (Chia Cheng) had had time to start a topic of conversation, the
        senior officer anticipated him, and speedily observed: "Your humble servant
        does not pay this visit to-day to your worthy mansion on his own authority,
        but entirely in compliance with instructions received, as there is a favour
        that I have to beg of you.}
}
{
}

\Verse{13}
{
    \zA{宝玉如得了珍宝，便赶上来拉他，说道：“快 进去告诉：老爷要打我呢!快去，快去!要紧，要紧！” 宝玉一则急了说话不明白，
        二则老婆子偏偏又耳聋，不曾听见是什么话，把“要紧”二字只听做“跳井”二字， 便笑道：“跳井让他跳去，二爷怕什么？” 宝玉见是个聋子，便着急道：“你出去
        叫我的小厮来罢！” 那婆子道：“有什么不了的事？ 老早的完了。 太太又赏了银子， 怎么不了事呢？”
        宝玉急的手脚正没抓寻处，只见贾政的小厮走来，逼着他出去了。}
}
{
    \eA{I make bold to trouble you, esteemed Sir, on behalf of his highness, to take
        any steps you might deem suitable, and if you do, not only will his highness
        remember your kindness, but even I, your humble servant, and my colleagues
        will feel extremely grateful to you." Chia Cheng listened to him, but he
        could not nevertheless get a clue of what he was driving at. Promptly
        returning his smile, he rose to his feet.}
}
{
}

\Verse{14}
{
    \zA{贾政一见， 眼都红了，也不暇问他在外流荡优伶，表赠私物，在家荒疏学业，逼淫母婢，只喝 命：“堵起嘴来，着实打死！”
        小厮们不敢违，只得将宝玉按在凳上，举起大板， 打了十来下。 宝玉自知不能讨饶，只是呜呜的哭。 贾政还嫌打的轻，一脚踢开掌板
        的，自己夺过板子来，狠命的又打了十几下。 宝玉生来未经过这样苦楚，起先觉得 打的疼不过还乱嚷乱哭，后来渐渐气弱声嘶，哽咽不出。 众门客见打的不祥了，赶
        着上来，恳求夺劝。 贾政那里肯听？}
}
{
    \eA{"You come, Sir," he inquired, "at the instance of his royal highness, but
        what, I wonder, are the commands you have to give me? I hope you will
        explain them to your humble servant, worthy Sir, in order to enable him to
        carry them out effectively." The senior officer gave a sardonic smile.
        "There's nothing to carry out," he said.}
}
{
}

\Verse{15}
{
    \zA{说道：“你们问问他干的勾当，可饶不可饶!素 日皆是你们这些人把他酿坏了，到这步田地，还来劝解!明日酿到他弑父弑君，你 们才不劝不成？”
        众人听这话不好，知道气急了，忙乱着觅人进去给信。 王夫人听 了，不及去回贾母，便忙穿衣出来，也不顾有人没人，忙忙扶了一个丫头赶往书房
        中来，慌得众门客小厮等避之不及。 贾政正要再打，一见王夫人进来，更加火上浇油，那板子越下去的又狠又快。 按宝玉的两个小厮忙松手走开，宝玉早已动弹不得了。}
}
{
    \eA{"All you, venerable Sir, have to do is to utter one single word and the
        whole thing will be effected. There is in our mansion a certain Ch'i Kuan,
        who plays the part of young ladies. He hitherto stayed quietly in the
        mansion; but for the last three or five days or so no one has seen him
        return home. Search has been instituted in every locality, yet his
        whereabouts cannot be discovered.}
}
{
}

\Verse{16}
{
    \zA{贾政还欲打时，早被王夫人 抱住板子。 贾政道：“罢了，罢了!今日必定要气死我才罢！” 王夫人哭道：“宝 玉虽然该打，老爷也要保重。
        且炎暑天气，老太太身上又不大好，打死宝玉事小， 倘或老太太一时不自在了，岂不事大？” 贾政冷笑道：“倒休提这话!我养了这不 肖的孽障，我已不孝；
        平昔教训他一番，又有众人护持。 不如趁今日结果了他的狗 命，以绝将来之患！” 说着，便要绳来勒死。 王夫人连忙抱住哭道：“老爷虽然应
        当管教儿子，也要看夫妻分上。}
}
{
    \eA{But throughout these various inquiries, eight out of the ten tenths of the
        inhabitants of the city have, with one consent, asserted that he has of late
        been on very friendly terms with that honourable son of yours, who was born
        with the jade in his mouth.}
}
{
}

\Verse{17}
{
    \zA{我如今已五十岁的人，只有这个孽障，必定苦苦的 以他为法，我也不敢深劝。 今日越发要弄死他，岂不是有意绝我呢？ 既要勒死他，
        索性先勒死我，再勒死他!我们娘儿们不如一同死了，在阴司里也得个倚靠。” 说 毕，抱住宝玉，放声大哭起来。 贾政听了此话，不觉长叹一声，向椅上坐了，泪如
        雨下。 王夫人抱着宝玉，只见他面白气弱，底下穿着一条绿纱小衣，一片皆是血渍。}
}
{
    \eA{This report was told your servant and his colleagues, but as your worthy
        mansion is unlike such residences as we can take upon ourselves to enter and
        search with impunity, we felt under the necessity of laying the matter
        before our imperial master. 'Had it been any of the other actors,' his
        highness also says, 'I wouldn't have minded if even one hundred of them had
        disappeared;}
}
{
}

\Verse{18}
{
    \zA{禁不住解下汗巾去，由腿看至臀胫，或青或紫，或整或破，竟无一点好处，不觉失 声大哭起“苦命的儿”来。 因哭出“苦命儿”来，又想起贾珠来，便叫着贾珠哭道：
        “若有你活着，便死一百个我也不管了！” 此时里面的人闻得王夫人出来，李纨、 凤姐及迎、探姊妹两个也都出来了。 王夫人哭着贾珠的名字，别人还可，惟有李纨
        禁不住也抽抽搭搭的哭起来了。 贾政听了，那泪更似走珠一般滚了下来。 正没开交处，忽听丫鬟来说：“老太太来了！”}
}
{
    \eA{but this Ch'i Kuan has always been so ready with pat repartee, so respectful
        and trustworthy that he has thoroughly won my aged heart, and I could never
        do without him.'}
}
{
}

\Verse{19}
{
    \zA{一言未了，只听窗外颤巍巍的 声气说道：“先打死我，再打死他，就干净了！” 贾政见母亲来了，又急又痛，连 忙迎出来。 只见贾母扶着丫头，摇头喘气的走来。
        贾政上前躬身陪笑说道：“大暑 热的天，老太太有什么吩咐，何必自己走来，只叫儿子进去吩咐便了。” 贾母听了，
        便止步喘息，一面厉声道：“你原来和我说话!我倒有话吩咐，只是我一生没养个 好儿子，却叫我和谁说去！” 贾政听这话不像，忙跪下含泪说道：“儿子管他，也
        为的是光宗耀祖。}
}
{
    \eA{He entreats you, therefore, worthy Sir, to, in your turn, plead with your
        illustrious scion, and request him to let Ch'i Kuan go back, in order that
        the feelings, which prompt the Prince to make such earnest supplications,
        may, in the first place, be satisfied: and that, in the next, your mean
        servant and his associates may be spared the fatigue of toiling and
        searching."}
}
{
}

\Verse{20}
{
    \zA{老太太这话，儿子如何当的起？” 贾母听说，便啐了一口，说道： “我说了一句话，你就禁不起!你那样下死手的板子，难道宝玉儿就禁的起了？ 你说
        教训儿子是光宗耀祖，当日你父亲怎么教训你来着。” 说着也不觉泪往下流。 贾政 又陪笑道：“老太太也不必伤感，都是儿子一时性急，从此以后再不打他了。” 贾
        母便冷笑两声道：“你也不必和我赌气，你的儿子，自然你要打就打。 想来你也厌 烦我们娘儿们，不如我们早离了你，大家干净。”}
}
{
    \eA{At the conclusion of this appeal, he promptly made a low bow. As soon as
        Chia Cheng found out the object of his errand, he felt both astonishment and
        displeasure. With all promptitude, he issued directions that Pao-yü should
        be told to come out of the garden. Pao-yü had no notion whatever why he was
        wanted. So speedily he hurried to appear before his father. "What a regular
        scoundrel you are!" Chia Cheng exclaimed.}
}
{
}

\Verse{21}
{
    \zA{说着，便令人：“去看轿!我和 你太太、宝玉儿立刻回南京去！” 家下人只得答应着。 贾母又叫王夫人道：“你也 不必哭了。 如今宝玉儿年纪小，你疼他；
        他将来长大，为官作宦的，也未必想着你 是他母亲了。 你如今倒是不疼他，只怕将来还少生一口气呢！” 贾政听说，忙叩头
        说道：“母亲如此说，儿子无立足之地了。” 贾母冷笑道：“你分明使我无立足之 地，你反说起你来!只是我们回去了，你心里干净，看有谁来不许你打！”}
}
{
    \eA{"It is enough that you won't read your books at home; but will you also go
        in for all these lawless and wrongful acts? That Ch'i Kuan is a person whose
        present honourable duties are to act as an attendant on his highness the
        Prince of Chung Shun, and how extremely heedless of propriety must you be to
        have enticed him, without good cause, to come away, and thus have now
        brought calamity upon me?"}
}
{
}

\Verse{22}
{
    \zA{一面说， 一面只命：“快打点行李车辆轿马回去！” 贾政直挺挺跪着，叩头谢罪。 贾母一面说，一面来看宝玉。 只见今日这顿打不比往日，又是心疼，又是生气，
        也抱着哭个不了。 王夫人与凤姐等解劝了一会，方渐渐的止住。 早有丫鬟媳妇等上 来要搀宝玉。
        凤姐便骂：“糊涂东西!也不睁开眼瞧瞧，这个样儿，怎么搀着走的？ 还不快进去把那藤屉子春凳抬出来呢！” 众人听了，连忙飞跑进去，果然抬出春凳
        来，将宝玉放上，随着贾母王夫人等进去，送至贾母屋里。}
}
{
    \eA{These reproaches plunged Pao-yü in a dreadful state of consternation. With
        alacrity he said by way of reply: "I really don't know anything about the
        matter! To what do, after all, the two words Ch'i Kuan refer, I wonder!
        Still less, besides, am I aware what entice can imply!" As he spoke, he
        started crying.}
}
{
}

\Verse{23}
{
    \zA{彼时贾政见贾母怒气未消，不敢自便，也跟着进来。 看看宝玉果然打重了，再 看看王夫人一声“肉”一声“儿”的哭道：“你替珠儿早死了，留着珠儿，也免你
        父亲生气，我也不白操这半世的心了!这会子你倘或有个好歹，撂下我，叫我靠那 一个？” 数落一场，又哭“不争气的儿”。 贾政听了，也就灰心自己不该下毒手打
        到如此地步。 先劝贾母，贾母含泪说道：“儿子不好，原是要管的，不该打到这个 分儿。}
}
{
    \eA{But before Chia Cheng could open his month to pass any further remarks,
        "Young gentleman," he heard the senior officer interpose with a sardonic
        smile: "you shouldn't conceal anything! if he be either hidden in your home,
        or if you know his whereabouts, divulge the truth at once; so that less
        trouble should fall to our lot than otherwise would. And will we not then
        bear in mind your virtue, worthy scion!"}
}
{
}

\Verse{24}
{
    \zA{你不出去，还在这里做什么!难道于心不足，还要眼看着他死了才算吗？” 贾政听说，方诺诺的退出去了。 此时薛姨妈、宝钗、香菱、袭人、湘云等也都在这里。
        袭人满心委屈，只不好 十分使出来。 见众人围着，灌水的灌水，打扇的打扇，自己插不下手去，便索性走
        出门，到二门前，命小厮们找了焙茗来细问：“方才好端端的，为什么打起来？ 你 也不早来透个信儿！” 焙茗急的说：“偏我没在跟前，打到半中间，我才听见了。}
}
{
    \eA{"I positively don't know." Pao-yü time after time maintained. "There must, I
        fear, be some false rumour abroad; for I haven't so much as seen anything of
        him." The senior officer gave two loud smiles, full of derision. "There's
        evidence at hand," he rejoined, "so if you compel me to speak out before
        your venerable father, won't you, young man, have to suffer the
        consequences?}
}
{
}

\Verse{25}
{
    \zA{忙打听原故，却是为琪官儿和金钏儿姐姐的事。” 袭人道：“老爷怎么知道了？” 焙茗道：“那琪官儿的事，多半是薛大爷素昔吃醋，没法儿出气，不知在外头挑唆
        了谁来，在老爷跟前下的蛆。 那金钏儿姐姐的事，大约是三爷说的，我也是听见跟 老爷的人说。” 袭人听了这两件事都对景，心中也就信了八九分。 然后回来，只见
        众人都替宝玉疗治。 调停完备，贾母命：“好生抬到他屋里去。” 众人一声答应， 七手八脚，忙把宝玉送入怡红院内自己床上卧好。}
}
{
    \eA{But as you assert that you don't know who this person is, how is it that
        that red sash has come to be attached to your waist?" When Pao-yü caught
        this allusion, he suddenly felt quite out of his senses. He stared and
        gaped; while within himself, he argued: "How has he come to hear anything
        about this! But since he knows all these secret particulars, I cannot, I
        expect, put him off in other points;}
}
{
}

\Verse{26}
{
    \zA{又乱了半日，众人渐渐的散去了。 袭人方才进前来，经心服侍细问。 要知端底，究竟如何，且听下回分解。 (本章完)}
}
{
    \eA{so wouldn't it be better for me to pack him off, in order to obviate his
        blubbering anything more?" "Sir," he consequently remarked aloud, "how is it
        that despite your acquaintance with all these minute details, you have no
        inkling of his having purchased a house? Are you ignorant of an essential
        point like this?}
}
{
}

\Verse{27}
{
    \zA{}
}
{
    \eA{I've heard people say that he's, at present, staying in the eastern suburbs
        at a distance of twenty li from the city walls; at some place or other
        called Tzu T'an Pao, and that he has bought there several acres of land and
        a few houses. So I presume he's to be found in that locality; but of course
        there's no saying." "According to your version," smiled the senior officer,
        as soon as he heard his explanation, "he must for a certainty be there. I
        shall therefore go and look for him. If he's there, well and good; but if
        not, I shall come again and request you to give me further directions."
        These words were still on his lips, when he took his leave and walked off
        with hurried step.}
}
{
}

\Verse{28}
{
    \zA{}
}
{
    \eA{Chia Cheng was by this time stirred up to such a pitch of indignation that
        his eyes stared aghast, and his mouth opened in bewilderment; and as he
        escorted the officer out, he turned his head and bade Pao-yü not budge. "I
        have," (he said), "to ask you something on my return." Straightway he then
        went to see the officer off. But just as he was turning back, he casually
        came across Chia Huan and several servant-boys running wildly about in a
        body. "Quick, bring him here to me!" shouted Chia Cheng to the young boys.
        "I want to beat him." Chia Huan, at the sight of his father, was so
        terrified that his bones mollified and his tendons grew weak, and, promptly
        lowering his head, he stood still." "What are you running about for?"}
}
{
}

\Verse{29}
{
    \zA{}
}
{
    \eA{Chia Cheng asked. "These menials of yours do not mind you, but go who knows
        where, and let you roam about like a wild horse! Where are the attendants
        who wait on you at school?" he cried. When Chia Huan saw his father in such
        a dreadful rage, he availed himself of the first opportunity to try and
        clear himself. "I wasn't running about just now" he said. "But as I was
        passing by the side of that well, I caught sight, for in that well a
        servant-girl was drowned, of a human head that large, a body that swollen,
        floating about in really a frightful way and I therefore hastily rushed
        past." Chia Cheng was thunderstruck by this disclosure. "There's been
        nothing up, so who has gone and jumped into the well?" he inquired.}
}
{
}

\Verse{30}
{
    \zA{}
}
{
    \eA{"Never has there been anything of the kind in my house before! Ever since
        the time of our ancestors, servants have invariably been treated with
        clemency and consideration. But I expect that I must of late have become
        remiss in my domestic affairs, and that the managers must have arrogated to
        themselves the right of domineering and so been the cause of bringing about
        such calamities as violent deaths and disregard of life. Were these things
        to reach the ears of people outside, what will become of the reputation of
        our seniors? Call Chia Lien and Lai Ta here!" he shouted.}
}
{
}

\Verse{31}
{
    \zA{}
}
{
    \eA{The servant-lads signified their obedience, with one voice. They were about
        to go and summon them, when Chia Huan hastened to press forward. Grasping
        the lapel of Chia Cheng's coat, and clinging to his knees, he knelt down.
        "Father, why need you be angry?" he said. "Excluding the people in Madame
        Wang's rooms, this occurrence is entirely unknown to any of the rest; and I
        have heard my mother mention…." At this point, he turned his head, and cast
        a glance in all four quarters. Chia Cheng guessed his meaning, and made a
        sign with his eyes. The young boys grasped his purpose and drew far back on
        either side. Chia Huan resumed his confidences in a low tone of voice.}
}
{
}

\Verse{32}
{
    \zA{}
}
{
    \eA{"My mother," he resumed, "told me that when brother Pao-yü was, the other
        day, in Madame Wang's apartments, he seized her servant-maid Chin Ch'uan-erh
        with the intent of dishonouring her. That as he failed to carry out his
        design, he gave her a thrashing, which so exasperated Chin Ch'uan-erh that
        she threw herself into the well and committed suicide…." Before however he
        could conclude his account, Chia Cheng had been incensed to such a degree
        that his face assumed the colour of silver paper. "Bring Pao-yü here," he
        cried. While uttering these orders, he walked into the study.}
}
{
}

\Verse{33}
{
    \zA{}
}
{
    \eA{"If any one does again to-day come to dissuade me," he vociferated, "I shall
        take this official hat, and sash, my home and private property and surrender
        everything at once to him to go and bestow them upon Pao-yü; for if I cannot
        escape blame (with a son like the one I have), I mean to shave this scanty
        trouble-laden hair about my temples and go in search of some unsullied place
        where I can spend the rest of my days alone! I shall thus also avoid the
        crime of heaping, above, insult upon my predecessors, and, below, of having
        given birth to such a rebellious son."}
}
{
}

\Verse{34}
{
    \zA{}
}
{
    \eA{At the sight of Chia Cheng in this exasperation, the family companions and
        attendants speedily realised that Pao-yü must once more be the cause of it,
        and the whole posse hastened to withdraw from the study, biting their
        fingers and putting their tongues out. Chia Cheng panted with excitement. He
        stretched his chest out and sat bolt upright on a chair. His whole face was
        covered with the traces of tears. "Bring Pao-yü! Bring Pao-yü!" he shouted
        consecutively. "Fetch a big stick; bring a rope and tie him up; close all
        the doors! If any one does communicate anything about it in the inner rooms,
        why, I'll immediately beat him to death." The servant-boys felt compelled to
        express their obedience with one consent, and some of them came to look
        after Pao-yü.}
}
{
}

\Verse{35}
{
    \zA{}
}
{
    \eA{As for Pao-yü, when he heard Chia Cheng enjoin him not to move, he forthwith
        became aware that the chances of an unpropitious issue outnumbered those of
        a propitious one, but how could he have had any idea that Chia Huan as well
        had put in his word? There he still stood in the pavilion, revolving in his
        mind how he could get some one to speed inside and deliver a message for
        him. But, as it happened, not a soul appeared. He was quite at a loss to
        know where even Pei Ming could be. His longing was at its height, when he
        perceived an old nurse come on the scene. The sight of her exulted Pao-yü,
        just as much as if he had obtained pearls or gems; and hurriedly approaching
        her, he dragged her and forced her to halt.}
}
{
}

\Verse{36}
{
    \zA{}
}
{
    \eA{"Go in," he urged, "at once and tell them that my father wishes to beat me
        to death. Be quick, be quick, for it's urgent, there's no time to be lost."
        But, first and foremost, Pao-yü's excitement was so intense that he spoke
        with indistinctness. In the second place, the old nurse was, as luck would
        have it, dull of hearing, so that she did not catch the drift of what he
        said, and she misconstrued the two words: "it's urgent," for the two
        representing jumped into the well. Readily smiling therefore: "If she wants
        to jump into the well, let her do so," she said. "What's there to make you
        fear, Master Secundus?"}
}
{
}

\Verse{37}
{
    \zA{}
}
{
    \eA{"Go out," pursued Pao-yü, in despair, on discovering that she was deaf, "and
        tell my page to come." "What's there left unsettled?" rejoined the old
        nurse. "Everything has been finished long ago! A tip has also been given
        them; so how is it things are not settled?" Pao-yü fidgetted with his hands
        and feet. He was just at his wits' ends, when he espied Chia Cheng's
        servant-boys come up and press him to go out. As soon as Chia Cheng caught
        sight of him, his eyes got quite red. Without even allowing himself any time
        to question him about his gadding about with actors, and the presents he
        gave them on the sly, during his absence from home; or about his playing the
        truant from school and lewdly importuning his mother's maid, during his stay
        at home, he simply shouted: "Gag his mouth and positively beat him till he
        dies!"}
}
{
}

\Verse{38}
{
    \zA{}
}
{
    \eA{The servant-boys did not have the boldness to disobey him. They were under
        the necessity of seizing Pao-yü, of stretching him on a bench, and of taking
        a heavy rattan and giving him about ten blows. Pao-yü knew well enough that
        he could not plead for mercy, and all he could do was to whimper and cry.
        Chia Cheng however found fault with the light blows they administered to
        him. With one kick he shoved the castigator aside, and snatching the rattan
        into his own hands, he spitefully let (Pao-yü) have ten blows and more.
        Pao-yü had not, from his very birth, experienced such anguish. From the
        outset, he found the pain unbearable; yet he could shout and weep as
        boisterously as ever he pleased;}
}
{
}

\Verse{39}
{
    \zA{}
}
{
    \eA{but so weak subsequently did his breath, little by little, become, so hoarse
        his voice, and so choked his throat that he could not bring out any sound.
        The family companions noticed that he was beaten in a way that might lead to
        an unpropitious end, and they drew near with all despatch and made earnest
        entreaties and exhortations. But would Chia Cheng listen to them? "You
        people," he answered, "had better ask him whether the tricks he has been up
        to deserve to be overlooked or not! It's you who have all along so
        thoroughly spoilt him as to make him reach this degree of depravity! And do
        you yet come to advise me to spare him? When by and bye you've incited him
        to commit parricide or regicide, you will at length, then, give up trying to
        dissuade me, eh?"}
}
{
}

\Verse{40}
{
    \zA{}
}
{
    \eA{This language jarred on the ears of the whole party; and knowing only too
        well that he was in an exasperated mood, they fussed about endeavouring to
        find some one to go in and convey the news. But Madame Wang did not presume
        to be the first to inform dowager lady Chia about it. Seeing no other course
        open to her, she hastily dressed herself and issued out of the garden.
        Without so much as worrying her mind as to whether there were any male
        inmates about or not, she straightway leant on a waiting-maid and hurriedly
        betook herself into the library, to the intense consternation of the
        companions, pages and all the men present, who could not manage to clear out
        of the way in time.}
}
{
}

\Verse{41}
{
    \zA{}
}
{
    \eA{Chia Cheng was on the point of further belabouring his son, when at the
        sight of Madame Wang walking in, his temper flared up with such increased
        violence, just as fire on which oil is poured, that the rod fell with
        greater spite and celerity. The two servant-boys, who held Pao-yü down,
        precipitately loosened their grip and beat a retreat. Pao-yü had long ago
        lost all power of movement. Chia Cheng, however, was again preparing to
        assail him, when the rattan was immediately locked tightly by Madame Wang,
        in both her arms. "Of course, of course," Chia Cheng exclaimed, "what you
        want to do to-day is to make me succumb to anger!" "Pao-yü does, I admit,
        merit to be beaten," sobbed Madame Wang; "but you should also, my lord, take
        good care of yourself!}
}
{
}

\Verse{42}
{
    \zA{}
}
{
    \eA{The weather, besides, is extremely hot, and our old lady is not feeling
        quite up to the mark. Were you to knock Pao-yü about and kill him, it would
        not matter much; but were perchance our venerable senior to suddenly fall
        ill, wouldn't it be a grave thing?" "Better not talk about such things!"
        observed Chia Cheng with a listless smile. "By my bringing up such a
        degenerate child of retribution I have myself become unfilial! Whenever I've
        had to call him to account, there has always been a whole crowd of you to
        screen him; so isn't it as well for me to avail myself of to-day to put an
        end to his cur-like existence and thus prevent future misfortune?" As he
        spoke, he asked for a rope to strangle him;}
}
{
}

\Verse{43}
{
    \zA{}
}
{
    \eA{but Madame Wang lost no time in clasping him in her embrace, and reasoning
        with him as she wept. "My lord and master," she said, "it is your duty, of
        course, to keep your son in proper order, but you should also regard the
        relationship of husband and wife. I'm already a woman of fifty and I've only
        got this scapegrace. Was there any need for you to give him such a bitter
        lesson? I wouldn't presume to use any strong dissuasion; but having, on this
        occasion, gone so far as to harbour the design of killing him, isn't this a
        fixed purpose on your part to cut short my own existence? But as you are
        bent upon strangling him, be quick and first strangle me before you strangle
        him!}
}
{
}

\Verse{44}
{
    \zA{}
}
{
    \eA{It will be as well that we, mother and son, should die together, so that if
        even we go to hell, we may be able to rely upon each other!" At the
        conclusion of these words, she enfolded Pao-yü in her embrace and raised her
        voice in loud sobs. After listening to her appeal, Chia Cheng could not
        restrain a deep sigh; and taking a seat on one of the chairs, the tears ran
        down his cheeks like drops of rain. But while Madame Wang held Pao-yü in her
        arms, she noticed that his face was sallow and his breath faint, and that
        his green gauze nether garments were all speckled with stains of blood, so
        she could not check her fingers from unloosening his girdle.}
}
{
}

\Verse{45}
{
    \zA{}
}
{
    \eA{And realising that from the thighs to the buttocks, his person was here
        green, there purple, here whole, there broken, and that there was, in fact,
        not the least bit, which had not sustained some injury, she of a sudden
        burst out in bitter lamentations for her offspring's wretched lot in life.
        But while bemoaning her unfortunate son, she again recalled to mind the
        memory of Chia Chu, and vehemently calling out "Chia Chu," she sobbed: "if
        but you were alive, I would not care if even one hundred died!"}
}
{
}

\Verse{46}
{
    \zA{}
}
{
    \eA{But by this time, the inmates of the inner rooms discovered that Madame Wang
        had gone out, and Li Kung-ts'ai, Wang Hsi-feng and Ting Ch'un and her
        sisters promptly rushed out of the garden and came to join her. While Madame
        Wang mentioned, with eyes bathed in tears, the name of Chia Chu, every one
        listened with composure, with the exception of Li Kung-ts'ai, who unable to
        curb her feelings also raised her voice in sobs. As soon as Chia Cheng heard
        her plaints, his tears trickled down with greater profusion, like pearls
        scattered about.}
}
{
}

\Verse{47}
{
    \zA{}
}
{
    \eA{But just as there seemed no prospect of their being consoled, a servant-girl
        was unawares heard to announce: "Our dowager lady has come!" Before this
        announcement was ended, her tremulous accents reached their ears from
        outside the window. "If you were to beat me to death and then despatch him,"
        she cried, "won't you be clear of us!" Chia Cheng, upon seeing that his
        mother was coming, felt distressed and pained. With all promptitude, he went
        out to meet her. He perceived his old parent, toddling along, leaning on the
        arm of a servant-girl, wagging her head and gasping for breath. Chia Cheng
        drew forward and made a curtsey. "On a hot broiling day like this," he
        ventured, forcing a smile, "what made you, mother, get so angry as to rush
        over in person?}
}
{
}

\Verse{48}
{
    \zA{}
}
{
    \eA{Had you anything to enjoin me, you could have sent for me, your son, and
        given me your orders." Old lady Chia, at these words, halted and panted.
        "Are you really chiding me?" she at the same time said in a stern tone.
        "It's I who should call you to task! But as the son, I've brought up, isn't
        worth a straw, to whom can I go and address a word?" When Chia Cheng heard
        language so unlike that generally used by her, he immediately fell on his
        knees. While doing all in his power to contain his tears: "The reason why,"
        he explained, "your son corrects his offspring is a desire to reflect lustre
        on his ancestors and splendour on his seniors; so how can I, your son,
        deserve the rebuke with which you greet me, mother?"}
}
{
}

\Verse{49}
{
    \zA{}
}
{
    \eA{At this reply, old lady Chia spurted contemptuously. "I made just one
        remark," she added, "and you couldn't stand it, and can Pao-yü likely put up
        with that death-working cane? You say that your object in correcting your
        son is to reflect lustre on your ancestors and splendour on your seniors,
        but in what manner did your father correct you in days gone by?" Saying
        this, tears suddenly rolled down from her eyes also. Chia Cheng forced
        another smile. "Mother;" he proceeded, "you shouldn't distress yourself!
        Your son did it in a sudden fit of rage, but from this time forth I won't
        touch him again." Dowager lady Chia smiled several loud sneering smiles.
        "But you shouldn't get into a huff with me!" she urged.}
}
{
}

\Verse{50}
{
    \zA{}
}
{
    \eA{"He's your son, so if you choose to flog him, you can naturally do so, but I
        cannot help thinking that you're sick and tired of me, your mother, of your
        wife and of your son, so wouldn't it be as well that we should get out of
        your way, the sooner the better, as we shall then be able to enjoy peace and
        quiet?" So speaking, "Go and look after the chairs." she speedily cried to a
        servant. "I and your lady as well as Pao-yü will, without delay, return to
        Nanking." The servant had no help but to assent. Old lady Chia thereupon
        called Madame Wang over to her. "You needn't indulge in sorrow!" she
        exhorted her. "Pao-yü is now young, and you cherish him fondly; but does it
        follow that when in years to come he becomes an official, he'll remember
        that you are his mother?}
}
{
}

\Verse{51}
{
    \zA{}
}
{
    \eA{You mustn't therefore at present lavish too much of your affection upon him,
        so that you may by and bye, spare yourself, at least, some displeasure."
        When these exhortations fell on Chia Cheng's ear, he instantly prostrated
        himself before her. "Your remarks mother," he observed, "cut the ground
        under your son's very feet." "You distinctly act in a way," cynically smiled
        old lady Chia, "sufficient to deprive me of any ground to stand upon, and
        then you, on the contrary, go and speak about yourself! But when we shall
        have gone back, your mind will be free of all trouble. We'll see then who'll
        interfere and dissuade you from beating people!"}
}
{
}

\Verse{52}
{
    \zA{}
}
{
    \eA{After this reply, she went on to give orders to directly get ready the
        baggage, carriages, chairs and horses necessary for their return. Chia Cheng
        stiffly and rigidly fell on his knees, and knocked his head before her, and
        pleaded guilty. Dowager lady Chia then addressed him some words, and as she
        did so, she came to have a look at Pao-yü. Upon perceiving that the
        thrashing he had got this time was unlike those of past occasions, she
        experienced both pain and resentment. So clasping him in her arms, she wept
        and wept incessantly. It was only after Madame Wang, lady Feng and the other
        ladies had reasoned with her for a time that they at length gradually
        succeeded in consoling her.}
}
{
}

\Verse{53}
{
    \zA{}
}
{
    \eA{But waiting-maids, married women, and other attendants soon came to support
        Pao-yü and take him away. Lady Feng however at once expostulated with them.
        "You stupid things," she exclaimed, won't you open your eyes and see! How
        ever could he be raised and made to walk in the state he's in! Don't you yet
        instantly run inside and fetch some rattan slings and a bench to carry him
        out of this on? At this suggestion, the servants rushed hurry-scurry inside
        and actually brought a bench; and, lifting Pao-yü, they placed him on it.
        Then following dowager lady Chia, Madame Wang and the other inmates into the
        inner part of the building, they carried him into his grandmother's
        apartments.}
}
{
}

\Verse{54}
{
    \zA{}
}
{
    \eA{But Chia Cheng did not fail to notice that his old mother's passion had not
        by this time yet abated, so without presuming to consult his own
        convenience, he too came inside after them. Here he discovered how heavily
        he had in reality castigated Pao-yü. Upon perceiving Madame Wang also
        crying, with one breath, "My flesh;" and, with another, saying with tears:
        "My son, if you had died sooner, instead of Chu Erh, and left Chu Erh behind
        you, you would have saved your father these fits of anger, and even I would
        not have had to fruitlessly worry and fret for half of my existence! Were
        anything to happen now to make you forsake me, upon whom will you have me
        depend?"}
}
{
}

\Verse{55}
{
    \zA{}
}
{
    \eA{And then after heaping reproaches upon herself for a time, break out afresh
        in lamentations for her, unavailing offspring, Chia Cheng was much cut up
        and felt conscious that he should not with his own hand have struck his son
        so ruthlessly as to bring him to this state, and he first and foremost
        directed his attention to consoling dowager lady Chia. "If your son isn't
        good," rejoined the old lady, repressing her tears, "it is naturally for you
        to exercise control over him. But you shouldn't beat him to such a pitch!
        Don't you yet bundle yourself away? What are you dallying in here for? Is it
        likely, pray, that your heart is not yet satisfied, and that you wish to
        feast your eyes by seeing him die before you go?"}
}
{
}

\Verse{56}
{
    \zA{}
}
{
    \eA{These taunts induced Chia Cheng to eventually withdraw out of the room. By
        this time, Mrs. Hsüeh together with Pao-ch'ai, Hsiang Ling, Hsi Jen, Shih
        Hsiang-yün and his other cousins had also congregated in the apartments. Hsi
        Jen's heart was overflowing with grief; but she could not very well give
        expression to it. When she saw that a whole company of people shut him in,
        some pouring water over him, others fanning him; and that she herself could
        not lend a hand in any way, she availed herself of a favourable moment to
        make her exit. Proceeding then as far as the second gate, she bade the
        servant-boys go and fetch Pei-Ming. On his arrival, she submitted him to a
        searching inquiry.}
}
{
}

\Verse{57}
{
    \zA{}
}
{
    \eA{"Why is it," she asked, "that he was beaten just now without the least
        provocation; and that you didn't run over soon to tell me a word about it?"
        "It happened," answered Pei Ming in great perplexity, "that I wasn't
        present. It was only after he had given him half the flogging that I heard
        what was going on, and lost no time in ascertaining what it was all about.
        It's on account of those affairs connected with Ch'i Kuan and that girl Chin
        Ch'uan." "How did these things come to master's knowledge?" inquired Hsi
        Jen. "As for that affair with Ch'i Kuan," continued Pei Ming, "it is very
        likely Mr. Hsüeh P'an who has let it out;}
}
{
}

\Verse{58}
{
    \zA{}
}
{
    \eA{for as he has ever been jealous, he may, in the absence of any other way of
        quenching his resentment, have instigated some one or other outside, who
        knows, to come and see master and add fuel to his anger. As for Chin
        Ch'uan-erh's affair it has presumably been told him by Master Tertius. This
        I heard from the lips of some person, who was in attendance upon master."
        Hsi Jen saw how much his two versions tallied with the true circumstances,
        so she readily credited the greater portion of what was told her.
        Subsequently, she returned inside. Here she found a whole crowd of people
        trying to do the best to benefit Pao-yü.}
}
{
}

\Verse{59}
{
    \zA{}
}
{
    \eA{But after they had completed every arrangement, dowager lady Chia impressed
        on their minds that it would be better were they to carefully move him into
        his own quarters. With one voice they all signified their approval, and with
        a good deal of bustling and fussing, they speedily transferred Pao-yü into
        the I Hung court, where they stretched him out comfortably on his own bed.
        Then after some further excitement, the members of the family began
        gradually to disperse. Hsi Jen at last entered his room, and waited upon him
        with singleness of heart. But, reader, if you feel any curiosity to hear
        what follows, listen to what you will find divulged in the next chapter.}
}
{
}
